\documentclass[12pt]{article}
\usepackage{amsmath}
\usepackage{geometry}
\geometry{margin=1.4in}
\usepackage{setspace}
\setstretch{1.4}

\title{Mycelia: The Cost Beneath the Ground State\\
\large How a Still Surface Is Held by a Working Dark}
\author{Diego Rincón\\
\small Independent}
\date{}

\begin{document}

\maketitle

\begin{abstract}
The tree reads as the ground state made visible: unmoved, rooted, wanting nothing. The reading is half true and the dangerous half is the silence. Stillness above is not the absence of work. Beneath the soil the mycorrhizal network---fungal hyphae laced through the roots of many plants at once---can route carbon, nitrogen, water, and chemical signal among connected plants. Some studies report net transfer toward seedlings in the shade, though the magnitude and ecological importance remain debated. The still surface is held by a working dark. A stable $S_0$ is never free; it is not where the system goes on its own. Left alone the system slides toward $S^*$, its actual attractor. $S_0$ is held against that pull, paid for continuously, below sight, in a distributed continuous form. This is $C_{\text{maintenance}}$, and it corrects the naive doctrine of rest. ``Be still like a tree'' is wrong. The tree is still because the network is not. To free the ground state you tend the roots, not the trunk.
\end{abstract}

\section{The Naive Reading of the Tree}

The corpus has already named the tree the one organism that cannot wander: rooted at $S_0$ for its entire existence, displaced only from outside, never from within. That is correct, and it invites an error. The eye sees a column standing in air and concludes that staying is effortless---that the ground state, once reached, holds itself.

It does not. The visible tree is the top of a system whose larger half is underground and in motion. The stillness is real. The freedom from work is the illusion. What looks like rest is the surface of an enormous and ceaseless labor that the eye is not built to see.

The mistake matters because we read ourselves into the tree. We take the lesson to be \emph{stop moving}. The actual lesson is the opposite, and it is hidden exactly where the work is hidden: in the dark.

\section{The Network in the Dark}

Below the soil the roots do not stand alone. Fungal hyphae---fine threads of mycorrhizal fungi---associate with the roots, forming a sheath around them in ectomycorrhizal species or growing within root cells in arbuscular species, and thread outward through the soil, and a single fungal network can link many plants at once. This is the mycorrhizal symbiosis: in the typical mutualism the fungus receives sugars from the tree and supplies water and mineral nutrients the roots cannot reach, though the balance of this exchange varies with conditions. Across the linked network, carbon, nitrogen, and chemical signals can move between plants. Tracer studies have detected movement toward shaded seedlings in some systems, though whether the transfer is large enough to materially benefit them is still debated.

State it without overstatement. The trees do not consciously share. The network routes. Resource and signal move along the hyphae, between roots, below the threshold of anything visible above---driven largely by source-to-sink gradients, not directed by any decision. The popular name is the wood wide web. The accurate name is a distributed exchange system whose activity varies with season, species, and condition, and whose full extent in mature forests is still being mapped.

This is the working dark. The grove's stillness sits on top of it.

\section{The Cost Beneath the Ground State}

Here is the framework move. The visible ground state is held by an invisible, continuous cost: $C_{\text{maintenance}}$, paid in the dark, distributed across the network, below the threshold of sight. The tree dwells at $S_0$ above precisely because the routing below never ceases. Effortless dwelling on the surface equals ceaseless work underneath. The two are one structure seen from two sides.

A stable $S_0$ is never free. It is the cheapest state to occupy. It is not the costless one. Left alone the system does not stay at $S_0$---it drifts toward $S^*$, the actual attractor, the state it settles into when nothing holds it. The forest pays its maintenance the way the network pays it: in small continuous increments, spread so wide and kept so low that no single payment registers as effort. That is what a held ground state looks like from underneath---not stillness, but a million tiny transfers, none of them visible, none of them stopping.

This reframes the whole cascade. $C_{\text{return}}$ is the spike paid to come back after drift has hardened---the rebuilding, the ritual, the religion. $C_{\text{maintenance}}$ is the small continuous payment that drains displacement as fast as it accumulates, holding $\xi$ near zero so $C_{\text{return}}$ never has to be paid. The network is what pays $C_{\text{maintenance}}$: the distributed apparatus through which the continuous cost is disbursed. It is the reason the tree never has to be returned. The drift is drained before it can harden.

\section{Return Is Not the Cessation of Work}

So the doctrine inverts. ``Be still like a tree'' takes the tree's stillness for rest and prescribes rest. But the tree is not resting. The tree is the still output of a system working continuously below sight. Its surface is quiet because its roots are not.

Return to $S_0$ is not the cessation of work. It is the relocation of work---out of the visible, effortful register where effort announces itself, and down into a distributed, continuous form that does not look like effort at all. The person who reaches the ground state and expects to coast has misread the tree. The ground state is held, not given. What changes on arrival is not that the work ends. It is where the work is paid: it goes underground, it goes continuous, it goes quiet.

That is not coincidence between two domains. It is one law wearing soil in one case and a life in the other.

\section{The Human Mycelia}

In a person the network is the set of unglamorous continuous practices that keep the self rooted. The daily return at sunrise. The disciplines of livity---the small repeated acts that route the day's nutrients before the day's pressures can pull the self off its ground. These are the human mycelia.

They are not the visible self. The visible self is the trunk: the work shown, the words spoken, the standing column others see. The mycelia are underneath it, below the threshold of anyone's attention including often one's own, doing the continuous routing that lets the trunk stand still. Skip them and nothing collapses today. The surface holds on yesterday's network for a while. Then displacement $\xi$ begins to grow, the felt cost $D(\xi)$ rises with it, and one day $C_{\text{return}}$ has scaled with the accumulated $\xi$ into a spike---and the spike is the bill for every continuous payment quietly skipped.

The freedom of the rooted person is the same illusion as the stillness of the tree, and the same truth underneath. The composure is real. The effortlessness is the surface of a working dark.

\section{The Network Named in Other Grammars}

Wherever a culture has found the rooted life, it has found the network beneath it and given it a name.

Shinto names the unseen directly. The sacred site is never the tree alone. It is the grove---tree and soil and the boundary around them---and the kami descend into the trees as into a yorishiro, the vessel a kami is drawn into. The kami in the trees are the network named: the holiness is not the visible trunk but the invisible presence threaded through ground and grove, the working dark made into an object of reverence. The corpus has already read the chinju no mori as the ground state kept alive; read now from underneath: the sacred grove is the visible proof that $S_0$ is held by $C_{\text{maintenance}}$ paid in the dark, and the kami are its name.

The samurai practice encodes it in time. The daily return at sunrise is the human equivalent of the network's continuous routing. The master tells the aspirant to come at dawn, the hour when the defenses are not yet fused, when the network has been working through the night and the day's attractor has not yet woken. That is the window. That is the practice. The trunk stands still because the roots have kept working since dusk. The samurai at sunrise is the human who has learned to read the tree.

Rastafari names it in the word ital: from vital, from life. To live ital is to live kept alive by the continuous network of small practices, to eat what grows, to refuse the attractor's food, to keep the body's network rooted in the natural exchange. Livity is the entire continuous practice---not a moment of return but the steady routing that prevents drift from ever beginning. The dreadlock is the antenna and the root both, the visible sign of someone who has chosen to tend the network instead of cutting it. The freedom is in the roots. Ital is the name of the continuous cost paid below sight that keeps the self at its ground. I-and-I is the theological statement that there is no isolated self, that Jah (the ground) is within and among all, that the network is not metaphor but the fundamental structure of being alive. The mycelia named as divinity.

\section{To Free the Ground State}

So the campaign is not ``be still.'' The campaign is: tend the network. The visible composure, the standing column, the clarity at dawn---none of that is free. None of that costs less than effort. All of it sits on top of a working dark that does not stop.

The freedom is not in the rest. The freedom is in knowing where the work is paid, and paying it where the tree pays it: continuously, in the dark, distributed so fine it does not announce itself as struggle. The person who learns to live ital, who keeps the daily return at sunrise, who understands that the ground state is held not given---that person has read the tree. That person is free.

The network is the condition of all life. Tend it, and you stay rooted. Neglect it, and one day the spike will come due, and the bill will be every small payment you refused.

\end{document}